\section{Planar convex bodies}

In this section, we build up a theory for planar convex bodies that will be used extensively later.
We do not claim originality of the results described here.

\subsection{Convex body}

\begin{definition}
A set $C \subseteq \RR^2$ is \emph{convex} 
if for any two points $x, y \in C$, any point
$\lambda x + (1 - \lambda) y$ with arbitrary $0 \leq \lambda \leq 1$ on the line segment connecting $x$ and $y$ also belongs to $C$.
A set $C \subseteq \RR^2$ is a \emph{convex body} if it is a nonempty convex set which is compact.
\end{definition}
While convex polygons and convex regions of smooth (or piecewise $C^1$) boundaries are examples, the definition also allows the boundary to behave more wildly as well. 
A main goal of this section is to analyze such cases inclusively.

Remark that for a convex body $C$, $p_C \colon \RR / 2 \pi \ZZ \ra \RR$ is the support function of $C$.
We can define the tangent lines of $C$ using this function.
\begin{definition}
For a convex body $C$ and arbitrary $t \in \RR / 2 \pi \ZZ$, define the \emph{tangent line} $l_C(t) = \{x \in \RR^2 \colon x \cdot u_t = p_C(t)\}$ with the normal vector $u_t$.
\end{definition}

\begin{definition}
By the convexity and compactness of $C$, the intersection $e_C(t) = C \cap l_C(t)$ at an angle $t \in \RR / 2 \pi \ZZ$ is either a proper line segment or a point. If it is a point, we can consider it as a degenerate line segment with same endpoints.
Let the point $v_C^+(t)$ be the endpoint of $e_C(t)$ located at the farthest in the direction of $u_{t + \pi / 2}$.
Likewise, let the point $v_C^-(t)$ be the endpoint of $e_C(t)$ located at the farthest in the direction of $u_{t - \pi / 2}$.
\end{definition}
Note that if $C$ is a convex polygon, the points $v_C^+(t), v_C^-(t)$ for all $t \in \RR / 2 \pi \ZZ$ will describe the vertices of the polygon. 

\begin{definition}
For any $t_1, t_2 \in S^1$ such that $t_2 = t_1 + \theta$ for nonzero $\theta \in (-\pi, \pi)$,
the lines $l_C(t_1)$ and $l_C(t_2)$ are not parallel so they intersect.
Let $v_C(t_1, t_2)$ be the intersection point of two lines.
\end{definition}
\begin{proposition}
For any $t_1, t_2 \in S^1$ such that $t_2 = t_1 + \theta$ for nonzero $\theta \in (-\pi, \pi)$, we have 
$$
v_C(t_1, t_2) = p_C(t_1) u_{t_1} + \frac{p_C(t_2) - p_C(t_1) \cos \theta}{\sin \theta} v_{t_1}.
$$
\end{proposition}
\begin{proof}
Because the point $v = v_C(t_1, t_2)$ is on the line $l_C(t_1)$,
we have $v = p_C(t_1) u_{t_1} + \alpha v_{t_1}$
for some constant $\alpha \in \RR$.
We use $v \cdot u_2 = p_C(t_2)$ to derive $\alpha$.
\begin{align*}
p_C(t_2) &= v \cdot u_2 \\
&= p_C(t_1) (u_{t_1} \cdot u_{t_2}) + \alpha (v_{t_1} \cdot u_{t_2}) \\
&= p_C(t_1) \cos \theta + \alpha \sin \theta
\end{align*}
This gives $\alpha = p_C(t_2) \csc \theta - p_C(t_1) \cot \theta$ and completes the proof.
\end{proof}

\begin{proposition}
For any point $w \in \RR^2$ outside $C$, $w$ is equal to $v_C(t_1, t_2)$	
\end{proposition}


\begin{definition}
For any function $f \colon \RR \ra \RR$, denote its left derivative as $\partial_+f$ and right derivative as $\partial_- f$.
\end{definition}
\begin{proposition}
For any convex body $C$ and angle $t \in \RR / 2 \pi \ZZ$,
the following equalities hold.
$$v_C^+(t) = \lim_{t' \ra t^+} v_C(t, t') = p_C(t)u_t + \partial_+p_C(t)v_t$$
$$v_C^-(t) = \lim_{t' \ra t^-} v_C(t, t') = p_C(t)u_t + \partial_-p_C(t)v_t.$$

Moreover, 
for any $t \in S^1$ and $0 < \theta < \pi$,
if we define the values
$$l_C^+(t, t+\theta) = \frac{p_C(t_2) - p_C(t_1) \cos \theta}{\sin \theta} - \partial_+p_C(t)$$
and
$$l_C^-(t, t-\theta) = -\frac{p_C(t_2) - p_C(t_1) \cos \theta}{\sin \theta} + \partial_-p_C(t)$$
then $l_C^+(t, t+\theta)$ is the distance between $v^+_C(t)$ and $v_C(t, t+\theta)$,
and $l_C^-(t, t-\theta)$ is the distance
between $v^-_C(t)$ and $v_C(t, t-\theta)$,
and the values $l_C^+(t, t+\theta)$ and 
$l_C^-(t, t-\theta)$ are 
\end{proposition}
\begin{proof}
We first observe that 
\end{proof}

Most tangent lines touch $C$ at a single vertex.
\begin{proposition}
For any planar convex body $C$, there is only a countable number of angles $t \in S^1$ with the edge $e_C(t)$ being a proper line segment.
So the tangent line $l_C(t)$ touches $C$ at the unique point $v^+_C(t) = v^-_C(t)$ for all but a countable number of values of $t \in S^1$.
\end{proposition}
\begin{proof}[Proof sketch]
Use the monotonicity: for any $0 < t < \pi$, $v_C^+(t)$ has strictly lesser $x$-coordinate than $v_C^-(t)$ unless $v_C^+(t) = v_C^-(t)$. Also for any $0 < t < t' < \pi$, $v_C^-(t')$ has lesser or equal $x$-coordinate than $v_C^+(t)$.
Now if $v_C^+(t) \neq v_C^-(t)$ for uncountably many $t$'s, the $x$-coordinate of $v_C^-(\pi / 2)$ descends from the $x$-coordinate of $v_C^+(0)$ `uncountably many times', contradicting that $C$ is bounded. 
\end{proof}

\begin{corollary}
\label{cor:func}
For any planar convex body $C$, its support function $p_C$ is both left-differentiable and right-differentiable, and has only a countable number of $t \in S^1$ where the left and right differentiation do not match.

Also, if $C$ lies inside the ball of radius $R \geq 0$ centered at origin, then the functions $p_C, \partial_+p_C, \partial_-p_C$ are all bounded by $R$.
\end{corollary}

\subsection{Surface area measure}

We shall construct the \emph{surface area measure} $\beta$ of an arbitrary convex body $C$. 
Informally speaking, for any closed interval $I = [a, b]$ of $S^1$, $\beta(I)$ will be the length of the curve on the boundary of $C$ from $v^-(t_1)$ to $v^+(t_2)$.
In particular, the measure $\beta(\{t\})$ of a singleton $\{t\}$ is the distance between $v^-(t)$ and $v^+(t)$.
Consequently, $\beta([a, b))$ is the distance between $v^-(a)$ and $v^-(b)$ on the boundary, and $\beta((a, b])$ is the distance between $v^+(a)$ and $v^+(b)$ on the boundary.
The construction of such a measure is already known for $n$-dimensional convex bodies \textcolor{red}{TODO: Put citations}. 
However, as we will need show some formulas involving $\beta$ that specifically works for dimension 2, we will construct the surface area measure for dimension 2 in a different way and show later that it matches with the original definition.

Assume for a moment that $p_C$ is twice continuously differentiable.
Then for any angle $t \in S^1$ the tangent line $l_C(t)$ touches $C$ at the unique point $v_C(t) = v_C^+(t) = v_C^-(t) = p_C(t)u_t + p'_C(t)v_t$ on the boundary of $C$.
In particular, $v_C(t)$ parametrizes the boundary points of $C$.
Note that $v'_C(t) = (p_C(t) + p''_C(t))v_t$,
so that the length of the curve that connects $v_C(a)$ and $v_C(b)$ on the boundary of $C$ is 
$$\int_a^b ||v'_C(t)||dt 
= \int_a^b (p_C(t) + p''_C(t)) dt.
= \int_a^b p_C(t) dt + p'_C(b) - p'_C(a)$$
Therefore, in this case, we can say that $p_C(t) + p''_C(t)$ is the density function of the surface are measure $\beta$ we are looking for.

Now let's get back to the general setting where $C$ is an arbitrary convex body and $p_C(t)$ may not differentiable.
From the equation above, we can suspect that this will hold.
$$
\beta((a, b]) = \int_a^b p_C(t) dt + \partial^+_t p_C(b) - \partial^-_t p_C(a)
$$
We will now use this guess to actually define $\beta$.

\begin{definition}
For an arbitrary convex body $C$,
define the function $F_C \colon \RR \ra \RR$ as 
$$
F_C(x) = \int_0^x p_C(t)dt + \partial^+_t p_C(x) - \partial^-_t p_C(0)
$$
so that in particular, 
$$
F_C(b) - F_C(a) = \int_a^b p_C(t)dt + \partial^+_t p_C(b) - \partial^-_t p_C(a)
$$
follows for arbitrary $a, b \in \RR$.
\end{definition}

\begin{theorem}
For any convex body $C$, $F_C$ is monotonically increasing.
For any $a \in \RR$ and $b \in \RR$ so that $0 < b - a < \pi$,
we have $F_C(b) - F_C(a) \geq 0$.
\end{theorem}
\begin{proof}
Key ideas: 
First, the value is invariant under
translations of $C$.
That is, adding or subtracting multiples of $\cos t$ and $\sin t$ does not affect the value of $F_C(x)$.

Second, now we can set origin to any point that we want.
So we move $C$ and assume that $v_C(a, b)$ is the origin.
A consequence is that $p_C(t) \leq 0$ for all $t \in [a, b]$.
Another consequence is that $l_C^+(a) = \partial^+_tp_C(b)$, or whatever, that it will give us a positive number as well. This concludes the proof.  
\end{proof}

\begin{theorem}
$F_C$ is right-continuous.	
\end{theorem}
\begin{proof}
\url{https://math.stackexchange.com/questions/1757987/how-to-prove-that-the-right-derivative-of-a-convex-function-is-right-continuous}	
\end{proof}

Now $F_C$ gives rise to a measure $\beta_C$ (standard real analysis result -- need to search references).

Goal: NOT to show that $\beta_C$ measures boundary perimeter directly.
We do not use that.
Instead, observe that it works linearly,
and that for the area we do have the formula that is expected.

Trapezodial region, bounded by (i) the segment connecting origin and $v_C^+(a)$, (ii) the segment connecting origin and $v_C^+(b)$, and $l_C(a)$, and $l_C(b)$.
The area of this thing will have a closed form.


